\documentclass[11pt]{article}
\usepackage[margin=25mm]{geometry}
\usepackage{fontspec}
\setmainfont{Latin Modern Roman}
\usepackage{amsmath,amssymb,mathtools}
\usepackage{unicode-math}
\setmathfont{Latin Modern Math}
\usepackage{xurl,hyperref}
\hypersetup{colorlinks=true,urlcolor=blue}
\setcounter{secnumdepth}{0}
\setlength{\parindent}{0pt}
\setlength{\parskip}{5pt}
\setlength{\emergencystretch}{3em}
\newcommand{\sref}[2]{\hyperlink{source:#1:#2}{[S#2]}}
\newcommand{\cataloguescope}[1]{\paragraph{Recorded target.}#1}
\begin{document}
\section{1-Factorization Conjecture}\label{statement}

\subsection{Definitions and mathematical statement}

\paragraph{Shared notation.}$\mathbb N=\{1,2,\ldots\}$, and $\mathbb Z,\mathbb Q,\mathbb R,\mathbb C$ denote the integers, rationals, reals and complex numbers. Any other conventions are specified locally.


A finite simple undirected graph is $k$-regular if each vertex has degree $k$. A 1-factor is a perfect matching: a set of pairwise vertex-disjoint edges meeting every vertex once. A 1-factorization is a partition of the edge set into perfect matchings. Define $b(n)=n$ for odd $n$ and $b(n)=n-1$ for even $n$.
\paragraph{Statement recorded in the supplied source.}
For every $n\in\mathbb N$, every integer $k\ge b(n)$, and every $k$-regular simple graph $G$ on $2n$ vertices, must $G$ have a 1-factorization into $k$ perfect matchings? Equivalently, is the threshold $k\ge2\lceil n/2\rceil-1$ sufficient for every order $2n$? \sref{statement}{1}\sref{statement}{2}
The cited theorem proves this assertion for all sufficiently large $n$; the catalogue question here retains its stated quantifier over every $n$.
\subsection{Short English statement}
Does every regular graph on twice n vertices decompose into perfect matchings when its degree is at least n, or at least n minus one when n is even?
\subsection{Sources}
\begin{description}
\item[\hypertarget{source:statement:1}{S1}] ULAM.AI and the UnsolvedMath Contributors, \emph{UnsolvedMath: A Curated Collection of Open Mathematics Problems}, catalogue entry. CC BY 4.0. \url{https://huggingface.co/datasets/ulamai/UnsolvedMath}.
\item[\hypertarget{source:statement:2}{S2}] Béla Csaba, Daniela Kühn, Allan Lo, Deryk Osthus and Andrew Treglown, \emph{Proof of the 1-factorization and Hamilton decomposition conjectures}, Memoirs of the American Mathematical Society 244 (2016); arXiv:1401.4159v2, Section 1.1, Theorem 1.1.1. \url{https://arxiv.org/abs/1401.4159v2}.
\end{description}
\label{end:statement}
\end{document}
